\scanpage{162}
\mspage{\unsure{77|74}}
Dans ces termes, les topos envisagés comme ``pièces de construction'' sont des $V_{d_*,d'_*}$ et $V^*_{d_*,d'_*}$.
Du pt de vue notation, j'ai envie d'éliminer des redondances, telles que $D_{d_0,d_0}\xrightarrow{\sim}D_{d_0}$ et $D_{d_0,d_1,d_1}\xrightarrow{\sim}D_{d_0,d_1}$ (les op. \inserted{simpl.} de dégénérescence sont des iso!) en exigeant $d_0>\cdots>d_r$ dans la notation $D_{d_*}=D_{d_0,\ldots,d_r}$ (ce qui implique qu'on se borne à des op. simpliciales qui correspondent à des inclusions de simplexes types -- donc à des composés d'op. face), \deleted{et le} et les redondances du type $D_{d,0}\xrightarrow{\sim}D_d$, plus gén.
\[
D_{d_0,\ldots,d_r,0}\xrightarrow{\sim}D_{d_0,\ldots,d_r}\qquad(\text{si }d_r>0)
\]
en exigeant que si $r\geq1$, alors $d_r>0$. Avec cette notation, le seul $D_{d_*}$ où figure l'indice $0$ est $D_0$. Il faut donc compléter implicitement les op. simpliciales entre les $D_{d_*}$ en y ajoutant les projections
\[
D_{d_*}\unsure{-|\longrightarrow}D_0=X
\]
i.e. en regardant les $D_{d_*}$ comme étant des espaces analytiques \uline{au dessus de $X$} (les op. simpliciales respectant cette structure).

\margin{Ces morphismes \unsure{correspondraient} formellement à l'inclusion du système \uuline{vide} $\varnothing\subset\{d_0>\cdots>d_r\}$}
\deleted{\unsure{indépendamment}}
\margin{161}
